% \documentclass[a4paper, 12pt]{article}
%% \usepackage[left=3.9cm,right=3.9cm,top=3cm,bottom=3cm]{geometry}

\documentclass[a4paper,10pt,twocolumn]{article}
\usepackage[left=2.6cm,right=2.6cm,top=2.55cm,bottom=2.55cm]{geometry}

\usepackage{xurl} 
\usepackage{cancel}
\begin{document}

\title{Towards woman professors at DoCS}
\date{}

\maketitle

% \textsuperscript{,}\footnote{Released 8 March, 2026} first
% }

\paragraph{Abstract} When the author\footnote{Laura Marie Feeney, Uppsala Univeristy.} first
visited Uppsala University in the early 2000's, she was slightly
surprised and mildly disappointed that there were no women professors
at the Division of Computer Systems.
%% \footnote{Laura Marie Feeney is an experienced researcher with a
%% strong international profile. She has been a student, engineer,
%% researcher and manager at startups, companies, labs and universities
%% in the US and Sweden, and a visiting scholar at universities in
%% Germany and Italy.}
Over the following twenty years, DoCS has \emph{failed to hire and
retain even one woman} who would go on to be promoted to professor at
UU.  The situation should no longer be considered ``mildly
disappointing''.  It should be seen as evidence of a problem,
especially in context of a recent employee survey reporting no gender
discrimination at DoCS.

This essay\footnote{\url{https://github.com/lmf-sonri/TWPaDoCS}, v26.03.08}
discusses some gender-related aspects of hiring and retaining women
researchers, as well as issues related to evaluating and improving the
gender environment.  While the focus is on DoCS, some discussion is
more broadly applicable in the IT department and beyond.  It is
written from the perspective of an experienced Computer Systems
researcher, rather than a specialist in gender studies or HR policy
and so closes with some personal reflections.\footnote{The author has
no formal background in Gender Studies.  The intent is to address the
situation at DoCS, not to rehash ``women in tech''.}

%% The title is a traditional formulation for a paper that provides
%% useful structure, but no definitive solution, for some problem.

\subsection*{Introduction}

Women are broadly underrepresented in computer science and engineering
(CSE). In Europe and North America, slightly more than 20\% of PhDs in
CSE are awarded to women.\footnote{Unsurprisingly, estimates vary
somewhat, but recent values seem to be $\sim$19-22\%.}  However, this
statistic masks considerable variation among CSE sub-fields.  In the
US and Canada, it is estimated that 12\%-16\% of PhDs in systems
topics (architecture, networking, OS, robotics, and
security/assurance) were awarded to women in 2012-2019.  By contrast,
29\%-50\% of PhDs in the fields of HCI, IS, and CSE education were
awarded to women.\footnote{The CRA Taulbee survey:
\url{https://cra.org/crn/2019/08/expanding-the-
 pipeline-gender-and-ethnic-differences-in-phd-specialty-areas} has
surveyed US and Canadian CSE departments since the
1970's.}\textsuperscript{,}\footnote{No European organization seems to
collect comparable statistics.  Since the overall proportion of women
in CSE is similar in the two regions, it is reasonable to assume that
the distribution among sub-fields is also
similar.}\textsuperscript{,}\footnote{Approximately 10\% of authors at
international systems conferences in 2017 had women's names
(\url{https://arxiv.org/abs/2201.01757}).  This seems roughly
consistent with the estimate above, since this sample reflects
long-term ``leaky-pipeline'' demographics, not just recent PhDs.}

Women are both historically and currently under-represented at DoCS,
even taking into account the particularly low proportion of women in
computer systems.  Crucially, there does not seem to be any trend
toward sustainable improvement.

Over the past twenty years, the number of women lecturers
(assistant/associate professor) active at DoCS has varied between zero
and two, mostly the former.  To date, no woman lecturer has remained
at DoCS long enough to be promoted to (full) professor.  A handful of
women have been with DoCS in other roles. However, these have been
either short-term (e.g.\,post-doc) positions or secondary affiliations
with minimal presence at UU.
  
During this period, DoCS has been one of the larger divisions in the
IT department.  At least ten men have been promoted to professor,
though some are no longer active at DoCS.  The corps of senior
lecturers has generally had between eight and thirteen men, including
adjuncts.

Women have defended around 15\% of DoCS PhD theses since 2000, but the
number of women defending has decreased slightly over time.\footnote{A
quick DIVA search showed that more women defended theses in 2000-2012
(9) than in 2013-2025 (6). The sample is small and may be incomplete
(e.g.\,joint/external supervision, gender inferred from name in a few
cases), but the data certainly do not indicate strong growth.}  DoCS
currently has over 20 PhD students, but only three or four women will
defend theses at DoCS into the early 2030's.


\subsection*{Hiring women}

Successful recruiting requires that an organization is both willing to
hire women and able to bring in highly-qualified
candidates.\footnote{This essay addresses only the second issue.}
Most of the factors that make a position attractive to potential
applicants apply equally to men and women.  But there are also
gender-specific factors that may affect an organization's ability to
attract highly-qualified women.  If DoCS lags its peer institutions in
some way, it may face both immediate and accumulating challenges in
recruiting, given the limited pool of women candidates.

\paragraph{Visible gender environment}
Some women consider the gender environment when making career
decisions: \emph{Are there other women in my field working there and
are they successful?}

Potential applicants have many sources of information about DoCS and
the IT department.  Departmental web pages identify women associated
with DoCS, while tools like Google Scholar have reliable information
about their research activity and current affiliations.  This makes it
fairly easy to determine the actual representation of women.

The research community is also a source of subjective impressions of
universities' gender environments: PhD students and researchers
publish and participate widely in international conferences and
collaborations.  People can thus observe the career trajectories of
both junior and senior women in their field.\footnote{Edith Ngai's
move from UU to HKU was mentioned to the author at several
conferences.}  Organizations such as N2Women also sponsor networking
and mentorship events at major conferences.\footnote{``Networking
Networking Women''; the author has been a mentor in this program.}

Locally, IT students at UU experience some aspects of DoCS gender
environment, such as the balance between men and women instructors.
This can affect young women's academic and career outcomes, possibly
including whether they choose to continue at UU for MSc or
PhD.\footnote{The gender studies literature suggests that exposure to
women instructors and mentors has a positive effect.}  Visitors,
whether exchange students or senior scholars, may also get a sense of
the gender environment.

Finally, candidates who advance to the interview stage for a position
can directly observe some aspects of the DoCS gender environment.

\paragraph{Competition in hiring} The challenges faced by applicants
for academic jobs are widely discussed.  But universities (except for
the most elite) must also compete globally to make themselves
attractive to promising candidates.  Since only a small fraction of
potential applicants in computer systems are women, gender-related
factors can affect this competition.

As noted above, some women may prefer organizations that they view as
having a good gender environment and successful women researchers.
Mentors may also take their own impressions into account in advising
women.  These factors may affect women's choices in applying for
positions, even with a positive feedback effect.

Most organizations must follow rules about gender equality in
hiring. The formal mechanisms can vary widely, in both their specific
goals and their practical impact.  For example, some policies try to
make the hiring process as mechanical as possible, while others
actively promote hiring from under-represented groups.  Informally,
some departments, or even individuals, may be proactive in identifying
and encouraging applications from talented women.  Such efforts mean
that highly-qualified women may be quickly taken ``off the market'' by
organizations that have prioritized attracting them.

\paragraph{Impact}
The factors above are arguably minor compared to gender-independent
factors, such as funding, that are key to attracting good applicants.
Nevertheless, they can have a large impact: Typically, a few top-tier
applicants form the ``short-list'' of candidates for a position.
Statistically, it is likely that at most one of those candidates is a
woman.\footnote{I.e.\,assuming both candidate quality and applications
are roughly i.i.d wrt gender.  Math left as an exercise.}  Hiring is
therefore sensitive to factors that influence whether very
highly-qualified women choose to apply for or accept a position.  If
that one (statistical) woman who would otherwise have been on the
short-list does not apply, it may effectively eliminate the
possibility of hiring a woman.\footnote{\label{foot:hire}Note that if
some very highly-qualified man chooses not to apply, the short-list is
still mostly men.}\textsuperscript{,}\footnote{In this case, a rule
like ``interview at least one woman'' is useless; the women most
likely to successfully compete for the position on merit did not
apply.}

\subsection*{Retaining and promoting women}

DoCS appears to have substantial gender disparities with respect to
both the types of positions that senior researchers hold and their
retention rate.


Both men and women have had various kinds of secondary affiliations at
DoCS.  While such arrangements can be valuable for both sides, these
positions do not typically lead to a professorship.  Over the past
twenty years, it has often been the case that a \emph{large majority}
of the women who were listed as affiliated with DoCS have had their
main position elsewhere and only a minimal connection to UU.  By
contrast, most men affiliated with DoCS have had their main position
there.

Likewise, both men and women have left DoCS.  While it is neither
appropriate nor helpful to focus here on individuals, the fact is that
\emph{no woman} who has been hired as a lecturer has remained at DoCS
long enough to be promoted to professor.  By contrast, a high
proportion of men have remained at UU and many have been
promoted.\footnote{As with recruiting, retention is more sensitive to
failure to retain women lecturers (footnote
\footnotesize{\ref{foot:hire}}).}



Given the long-standing under-representation of women at DoCS, this
repeated failure to retain and promote women lecturers is cause for
concern: \emph{``Once is an accident, twice is a coincidence, three
times is enemy action.''}\footnote{An English saying. Statistics are
considered later.}

%% \textsuperscript{,}\footnote{There are serveral
%% versions.}\textsuperscript{,}\footnote{Why, yes, the author \emph{is} from Chicago.}

The decision to leave a position usually reflects some combination of
professional, personal and family considerations.\footnote{Also legal
issues like immigration, pension, etc.}  It is unsurprising that this
complex optimization problem sometimes resolves in favor of moving,
for both men and women.  Nevertheless, the gender environment in the
broader research community and in the workplace can affect women's
professional success and personal job satisfaction.  These factors can
influence a decision to seek new opportunities in academia or
industry.

\paragraph{Research community}
To be successful, a researcher needs to not only obtain high-quality
technical results, but also to establish themselves and their results
in the research community.  Competition for funding and attention is
therefore extremely intense -- high status and visibility in
professional networks is essential for attracting both.  Strong
relationships with key industry players are also crucial for securing
certain kinds of funding and for driving real-world impact.

%% for access to top-tier
%% collaborations and funding, as well as for CV-enhancers like invited
%% talks and prestigious service roles. 

But as researchers advance in their careers, perceptions and
evaluations of their professional merit necessarily become more
subjective: The ``right answer'' is not known and novel research is
often initially uncertain or incomplete.  Experts must rely in part on
their intuition and judgment in assessing which research directions
and results -- and which researchers -- are most convincing and
promising.

%% Breaking new ground technically is a complex and uncertain venture.
%% There are uncertain results and false starts along the way, and no
%% clear figure of merit.

Technical expertise is not enough; professional success depends in
part on social competition and human reactions.  These can be affected
by gender dynamics, especially in fields like computer systems, where
women are severely under-represented and sometimes perceived as
receiving preferential treatment.  For example, credit for joint
achievements may not be shared or perceived fairly.  Even ostensibly
objective and gender-neutral measures -- publication metrics, teaching
evaluations, collaborations, funding, and awards -- partly reflect the
cumulative impact of subjective social factors.

To thrive as researchers, women must be able to operate on a fully
equal basis, not only within IT department, but also in the Swedish
and international research communities.  Although the latter are
obviously outside the university's direct control, the university and
department should be institutionally proactive about this.  It should
not be incumbent on individuals to monitor and advocate for gender
equality, especially at the governmental level.

% if UU aspires to a leading role in Swedish and European society,

\paragraph{Local gender environment}

In the recent five-yearly survey of IT Department, DoCS employees
reported no gender discrimination in DoCS.\footnote{Only $\sim$1/3 of
employees responded; it obviously cannot be known whether any women
did.  (The author did not.)}  Nevertheless, both DoCS' long-term
failure to retain women lecturers -- and the survey result itself --
suggest that there are grounds for concern about the gender
environment at DoCS.

Historically, there have been at most one or two women substantially
active at DoCS in senior roles. In the gender studies literature,
women who have few or no other women as colleagues, role models or
mentors are reported to have lower job satisfaction and lower
retention rates.  Suggested causes include poor organizational
culture, feelings of isolation, pressure to act in a representative
role and a sense of facing a ``non-level playing field''.\footnote{The
author has limited competency with the gender studies literature, but
levels of significance seem to vary.\label{disclaimer}}

As noted earlier, a number of women have had a secondary affiliation
with DoCS; often these women have made up majority of women in the
division.  This practice may provide a little more presence for women
at DoCS (and improve its nominal gender statistics).  But it has also
created an unfortunate pattern where women ``at'' DoCS typically have
little activity or influence there.

There is a certain amount of bureaucracy and cosmetic norms around
gender equality in the IT department and university.  To some extent,
it is treated as (not entirely incorrectly) as meaningless process in
an organization that is already awash with it.  In any case, it has
resulted in a strange equilibrium between procedural ``success'' and
actual failure to achieve substantial improvement in DoCS' gender
environment.

\subsection*{Evaluating the gender environment}

Although DoCS is one of the larger divisions in the IT department, it
is not a statistically large sample, especially given the low
proportion of women in computer systems.  It is therefore important to
understand the role and limitations of any evaluation method, both
with respect to DoCS and more generally.

Statistical checks can compare easily controlled values, such as
salary or teaching hours.  As noted earlier, however, these simple
metrics are not the ones that have the most impact on a researcher's
career.  In theory, statistics can also address somewhat pointless
hypotheticals like \emph{What is the probability that DoCS' current
and historical gender statistics can be explained purely by
chance?}\footnote{I.e.\ compared to the gender statistics obtained via
a random process of arrivals and departures.}

% \footnote{For example, women instructors are often seen as
% less knowledgeable in controlled experiments.} 

Given the sample size constraints it may be impossible to make strong
statistical claims about gender equality in DoCS.  Moreover, there are
significant limitations (discussed below) to statistical and survey
methods that reduce the effectiveness of these tools in gender
unbalanced environments.  Both of these factors support the natural
bureaucratic tendency to focus on ``paperwork'' compliance, rather
than grapple with sensitive problems.

\paragraph*{Methodological aside (technical)} Statistical and survey methods are
widely used for measuring gender equality, both for evaluating
specific situations and for gender studies research on ``women in
STEM''.  Many studies make an implicit assumption that the sample is
fairly homogeneous.  However, the proportion of women varies widely
both between fields (e.g.\,biology \emph{vs} CSE) and within fields
(e.g.\,HCI \emph{vs} computer systems).  Samples that span fields or
sub-fields may therefore reflect a mix of gender environments:
Potential examples include university-, department-, or agency-level
studies, as well as work in gender research, where large $N$ samples
are often desirable.

In such cases, most of the women in the sample will necessarily come
from precisely those environments where women are most \emph{strongly
represented}.  The results will therefore be dominated by their
experiences, rather than those of the much smaller set of women from
environments where they are most \emph{severely
under-represented}.\footnote{Consider an (numerically convenient, but
not wildly unrealistic) hypothetical CSE department with 20 women and
80 men.  Division A has 9 women and 14 men and division B has 8 women
and 12 men, while divisions C, D, and E each have 1 woman and 18 men.
Any department-wide measurement or survey overwhelmingly (85\%)
reflects the experiences of the women in the two well-balanced
divisions A and B.  But in divisions C, D, and E -- a majority (57\%)
of the department -- women are outnumbered 18 to
1.}\textsuperscript{,}\footnote{ Uncontrolled sampling -- ``please
forward this survey to your female colleagues'' -- can amplify this
effect.  Consider an (over-)simplified example where a survey is sent
to a woman in a field where women have (on average) eight women
colleagues and to a woman in a field where they have only three ($<$3x
difference).  After two rounds of forwarding, the former group has had
at most 57 unique contacts with potential respondents and the latter
has had at most 7 ($>$8x difference).  (Somebody (else) should model
this properly, perhaps on co-author networks.)}

It is clearly the latter group that demands particularly careful
scrutiny with respect to gender equality.  But these women are nearly
invisible in the aggregated data!  Such studies may therefore be less
complete and accurate than assumed.

Privacy is also an issue when investigating gender issues in very
unbalanced environments.  Even ``anonymous'' results provide little or
no privacy to the women involved, since any gender-related content or
analysis necessarily reflects the data associated with a very small
number of women.\footnote{In our hypothetical CSE department,
division-level gender comparisons are anonymous for both men and women
in divisions A and B, but reveal the data associated with the lone
women in divisions C, D and E.}\textsuperscript{,}\footnote{Recall
that no gender discrimination was reported in the DoCS employee
survey.}

Despite the methodological challenges, it is important -- both for
monitoring gender equality within organizations and for gender studies
research -- to find ways to effectively measure environments where
there are very few or even zero women.\footnote{You could call it the
``dark matter'' of gender studies.}
%% \textsuperscript{,}\footnote{If you have ``physics envy''.}
%% \textsuperscript{,}\footnote{Which is different from Freudian
%% ``penis envy''.}\textsuperscript{,}\footnote{Nah, not really.}

\subsection*{Conclusion and personal reflections}

In the author's opinion, there are long-standing issues with the
gender environment at DoCS.  A disproportionate failure to hire, and
especially to retain, women lecturers is particularly worrisome.  Over
the past 20 years, DoCS has made almost no sustained progress in
increasing the number of senior women researchers active in the
division.  Nor has it significantly grown the pipeline of women at the
PhD and post-doc levels.

Moreover, the recent survey of DoCS employees suggests some denial of
the problem (as well as highlighting the methodological limitations of
surveys discussed above).  This long-term failure may reduce women's
confidence in the seriousness of HR, IT and DoCS' leadership in
ensuring a good gender environment.  Weakness in the visible gender
environment may also reduce DoCS' international competitiveness.

It is naturally tempting to view a sensitive and challenging problem
as simply unfortunate or to argue that there is no proof that DoCS has
a notably poor record.\footnote{To be fair, the problem is hardly
limited to DoCS or even the IT department.}  Given the sample size,
this claim may even be statistically tenable.

%% Many organizations are working to improve their
%% gender environment, whether on their own imitative or by external
%% mandate.  DoCS needs to create a gender environment that allows it to
%% successfully recruit and retain highly-qualified women (fixme).

But achieving a good gender environment is an \emph{active goal}, not
a passive statistical outcome.  If DoCS were to have a comparable
pattern of long-term failure in applying for funding -- another
competitive process with many subjective and random external factors
-- it would probably not be viewed as just an unfortunate fluke.  The
correct response to twenty years of ``bad luck'' is a strategy for
getting some ``good luck''.\footnote{Raiding defenseless coastal
\cancel{villages} universities in viking longships, looting their
\cancel{gold} research grants, and stealing their women is currently
\cancel{encouraged} forbidden by the local authorities.
Inconvenient.}

The administrative incentives are overwhelmingly to \emph{document
that there is no problem}, not to address the problem.  Tick-box rules
and cosmetic norms such ``\emph{xyz} committee must have a woman'' may
do little more than put a disproportionally high administrative burden
on a small number of women researchers.

Moreover, emphasis on form contrasts unfavorably with persistent
failure on substance -- to create a gender environment that allows
DoCS to \emph{attract and retain highly-qualified women doing research
in computer systems}.  The gender studies literature seems to suggest
that this contradiction can create a kind of cognitive dissonance, or
even a sense that it is women who have failed due to lack of desire or
skill.\footnote{Students presumably observe the contradiction, too.
The take-home message may have as much impact on young men as on young
women.}\textsuperscript{,}\footnote{See also footnote
\footnotesize{\ref{disclaimer}}.}

%% \footnote{Almost as annoying as ``The funding application must
%% have a gender equality section.''}

It is wrong to dismiss criticism of DoCS' gender environment, even in
the absence of clear evidence of a specific policy violation.  Such an
attitude easily deteriorates into unwillingness and inability to work
toward meaningful evaluation and needed improvement.  Documenting
basic compliance with anti-discrimination rules is necessary.  But it
has clearly not been sufficient to create an environment that allows
DoCS to attract, retain and promote women researchers.  The former
should not be used to intimidate or suppress concerns about the
latter.

Finally, it is not the special responsibility of women at DoCS to fix
the issue.  Nor should they be required to provide ``constructive
solutions'', should they choose to raise it.  Instead, it is the
responsibility of the IT department, DoCS and the HR administration to
work seriously toward genuine improvement.

%% The gender environment at DoCS cannot be substantially improved
%% through mainly bureaucratic stratagems.  To increase the nominal
%% presence of women by adding external researchers or by attaching
%% researchers working far outside computer systems may do little more
%% than highlight that women Computer Systems.


\paragraph*{Personal reflections (appendix)}

Women in fields like computer systems (including the author) want to
focus on their work and succeed on their merits.  For the most part,
they just accept that their chosen area is overwhelmingly dominated by
men and get on with life.\footnote{More accurately, \emph{the women
who remain} are those who have been both willing and able to adapt to
this. Anyone with the ability to excel in computer systems has almost
certainly had many opportunities to choose to excel in something else
instead.}

Gender issues are mainly a distraction.  Life is hardly a controlled
experiment: It is rarely clear how talent, character and chance -- as
well as gender -- have contributed to success and failure.

In any case, the only sensible option is to succeed in the world as it
is.  Worrying about gender issues contributes nothing to that goal and
is also boring and annoying.  At worst, it risks becoming an easy
alternative to necessary self-criticism and diverting energy from
useful work and productive self-improvement.  It is more efficient to
focus on factors one can control.

But an unhealthy gender environment can have a broader impact than
occasionally irking the author and that creates a broader
responsibility.  In contexts where there are very few women, they
\emph{unavoidably} become representatives and role models, willing or
not.\footnote{Paraphrased from Ginni Rometty, ex-CEO of IBM.}  For the
author, this imposes a obligation to provide the strongest possible
representation and the responsibility grows with visibility and
seniority.

Strong representation obviously includes demonstrating excellence in
research -- pursuing an ambitious research agenda and obtaining
unassailably solid results.  A robust record of achievement in
research (and related KPI's) can serve as both a deterrent to
discrimination and a role model for women's ability to succeed.

%% Being in a
%% representative role can sometimes even create extra motivation (and
%% more pressure) in taking on challenges.

A duty to provide strong representation can be interpreted as
demanding nothing more than doing good research and providing good
mentorship.  That is what any serious researcher does anyhow, making
it a comfortable and sensible default: Research is already hard work;
gender issues are a distraction; and life is too short to waste on
stupid problems.

But years of good representation from the small number of women
researchers at DoCS has not led to any sustained improvement.
Fundamentally, the problem of an environment that consistently fails
to attract and retain women in senior roles \emph{has very little to
do with} ``excellence in research''.  

For the author, it is not clear how to ensure that women researchers
thrive at DoCS.  Nor is it clear how to balance personal pragmatism
with the obligations that come from being visible in a representative
position.  

No matter how tempting, equating representation with research
achievement risks irresponsibly sending women a message that the
situation is normal and acceptable or that they just need to ``excel
more'' to succeed.

Part of the difficulty is that any choice a woman makes --
nag\footnote{The author has very mixed feelings about ``women in
tech'' activism, some of which seems more performative than sensibly
grounded.}, focus on research, leave -- affects the gender
environment.  It is easy to inadvertently allow oneself to be used as
a ``token'' or to cover for organizational problems, which can
contribute to them being dismissed.  The obligation to provide strong
representation is not to an organization, but to women and society.

%% \textsuperscript{,}\footnote{If the author were to have used a word
%% similar to \emph{nonsense} at some point, she would have had to
%% delete it as more performative than sensibly grounded.}

Finding a balance is challenging.  It is hard to avoid simultaneous
criticisms of allowing oneself to be exploited because of gender and
of focusing on gender at the expense of research.  Perhaps there is no
peaceful solution, so a good researcher just writes up an intermediate
result and moves on.

%% %% \footnote{Equal opporutnity means that women should not need to be
%% %% exceptional to be equal.}

%% The obligation to ensure that women researchers thrive at DoCS is more
%% complex.  It is not clear what someone who is (however involuntarily)
%% in a representative position should do.  It is unwise to allow oneself
%% to be used as a ``token'' or to cover for organizational problems; the
%% core obligation is to women \emph{at} an organization.

%% % That the issue was seen as a non-problem in the employee survey may
%% % be partly the fault of the author.

%% Focusing on research is easy.  But when the gender environment is not
%% good, it is irresponsible to risk inadvertently sending the message
%% that the situation is acceptable or that women just need to ``excel
%% more''.  At the same time, any professional setback creates the
%% additional pressure of being a negative example to and about women.

%% %% The principle of equal opportunity means that women should not need to
%% %% be better to be equal.

%% There is an extremely difficult balance to achieve.  It is hard to
%% avoid simultaneous criticism of allowing oneself to be exploited and
%% of wasting energy on pointless problem.  Perhaps there is no peaceful
%% solution -- so a good researcher simply writes up an intermediate
%% result and moves on.

\paragraph{Acknowledgments}
The author appreciates the generous advice she has received about this
essay.

\end{document}


